\chapter{فضاءات باناخ والمبرهنات الأساسية}
\label{ch:b3:banach}

يدرس التحليل الدالي الفضاءات المعيارية اللامنتهية البُعد عبر
المؤثرات والمؤثرات الدالية التي تعيش عليها. واكتشافه المؤسِّس هو
أن \emph{التمام}، عبر مبرهنة بير، يفرض انتظامًا قويًا: فالعائلات
المحدودة نقطةً نقطة من المؤثرات محدودةٌ بالمعيار
(باناخ--شتاينهاوس)، وللتقابلات المتصلة مقلوباتٌ متصلة (التطبيق
المفتوح)، وتكشف المنحنيات البيانية الاتصالَ (المنحني البياني
المغلق). أما الركن الآخر، \emph{هان--باناخ}، فلا يحتاج إلى
التمام البتة --- بل إلى مبرهنة زورن المساعدة فحسب --- ويضمن أن
الفضاءات الثنوية غنية بما يكفي لرؤية كل متجهة. ويبرهن هذا الفصل
على المبرهنات الأربع كلها ويختبرها على فضاءات المتتاليات
الكلاسيكية $\ell^p$، وعلى حسابات ثنوية ملموسة، وعلى تطبيق مفاجئ
حقًّا: إذ توجد دوال متصلة دورية بالدور $2\pi$ \emph{تتباعد}
متسلسلة فورييه لها عند نقطة --- فيُحسم بالنفي سؤالٌ تركته السنة
الجامعية 2 مفتوحًا.

في كل ما يلي، $E, F$ فضاءان معياريان على $K = \R$ أو $\C$؛
وتعني عبارة \emph{فضاء باناخ}\index{فضاء باناخ} فضاءً معياريًا
تامًّا.

\section{المؤثرات المحدودة؛ فضاءات المتتاليات}

\begin{definition}\label{def:b3:banach:operator}
يرمز $\mathcal L(E, F)$ إلى فضاء التطبيقات الخطية \emph{المحدودة} (= المتصلة،
السنة الجامعية 2) وعليه \emph{معيار المؤثر}\index{معيار مؤثر}
$\vertiii T = \sup_{\norm x \leq
1}\norm{Tx}$؛ وهو تحت الضربي: $\vertiii{ST} \leq
\vertiii S\,\vertiii T$. و\emph{الفضاء
الثنوي}\index{فضاء ثنوي} هو $E' = \mathcal L(E, K)$.
\end{definition}

\begin{proposition}\label{prop:b3:banach:llcomplete}
إذا كان $F$ فضاء باناخ، فكذلك $\mathcal L(E, F)$؛ وعلى وجه الخصوص يكون $E'$
فضاء باناخ دائمًا.
\end{proposition}

\begin{proof}
لتكن $(T_n)$ متتالية كوشي من أجل $\vertiii\cdot$. فمن أجل كل
$x$، $\norm{T_nx - T_mx} \leq \vertiii{T_n - T_m}\norm x$: أي إن $(T_nx)$ متتالية كوشي في $F$، فهي متقاربة؛
ولنسمِّ نهايتها $Tx$. والتطبيق $T$ خطي (بنهايات المتطابقات
الخطية)؛ وبالانتقال إلى النهاية في $\norm{T_nx - T_mx} \leq \varepsilon\norm x$ (حيث $n, m \geq N$)
نجد $\norm{T_nx - Tx} \leq \varepsilon\norm x$: أي $T_n \to T$ \omterm{def:b3:banach:operator}{بمعيار المؤثر}، و $\vertiii T \leq \vertiii{T_N} +
\varepsilon < \infty$.
\end{proof}

\begin{definition}\label{def:b3:banach:lp}
فضاءات المتتاليات الكلاسيكية\index{فضاءات المتتاليات $\ell^p$}
(على $K$، مفهرسة بالدليل $\N$):
\[
\ell^p = \Bigl\{x = (x_n) : \norm x_p = \Bigl(\sum_n
\abs{x_n}^p\Bigr)^{1/p} < \infty\Bigr\}\quad (1 \leq p <
\infty),
\]
\[
\ell^\infty = \{x : \norm x_\infty = \sup\abs{x_n} < \infty\},
\]
و $c_0 = \{x : x_n \to 0\}$ مع $\norm\cdot_\infty$. وكون $\norm\cdot_p$ معيارًا ينتج من متراجحة
منكوفسكي، المبرهَن عليها في الحالة المتقطّعة تمامًا كما في
\cref{ch:b3:lp} (أو بجمع المتراجحة المنتهية البُعد من السنة الجامعية
2). وكلها فضاءات باناخ، ويكون $c_0$ فضاءً جزئيًا مغلقًا من
$\ell^\infty$ (\cref{exo:b3:banach:3}).
\end{definition}

\begin{omfigure}
\begin{tikzpicture}[scale=1.5]
  \draw[->] (-1.45,0) -- (1.45,0) node[right,font=\small]
    {$x_1$};
  \draw[->] (0,-1.45) -- (0,1.45) node[above,font=\small]
    {$x_2$};
  % p = 1 diamond
  \draw[omDef, thick] (1,0) -- (0,1) -- (-1,0) -- (0,-1)
    -- cycle;
  % p = 2 circle
  \draw[omProp, thick] (0,0) circle (1);
  % p = 4 superellipse |x|^4 + |y|^4 = 1
  \draw[omThm, thick, smooth]
    plot[domain=0:90, samples=46]
    ({cos(\x)/(cos(\x)^4 + sin(\x)^4)^0.25},
     {sin(\x)/(cos(\x)^4 + sin(\x)^4)^0.25});
  \draw[omThm, thick, smooth]
    plot[domain=90:180, samples=46]
    ({cos(\x)/(cos(\x)^4 + sin(\x)^4)^0.25},
     {sin(\x)/(cos(\x)^4 + sin(\x)^4)^0.25});
  \draw[omThm, thick, smooth]
    plot[domain=180:270, samples=46]
    ({cos(\x)/(cos(\x)^4 + sin(\x)^4)^0.25},
     {sin(\x)/(cos(\x)^4 + sin(\x)^4)^0.25});
  \draw[omThm, thick, smooth]
    plot[domain=270:360, samples=46]
    ({cos(\x)/(cos(\x)^4 + sin(\x)^4)^0.25},
     {sin(\x)/(cos(\x)^4 + sin(\x)^4)^0.25});
  % p = infty square
  \draw[black!60, thick] (-1,-1) rectangle (1,1);
  \node[omDef, font=\small] at (0.40,0.40) {$p{=}1$};
  \node[omProp, font=\small] at (-0.52,0.75) {$p{=}2$};
  \node[omThm, font=\small] at (0.90,-0.52) {$p{=}4$};
  \node[black!60, font=\small] at (-1.62,1.15)
    {$p{=}\infty$};
  \draw[black!60, ->] (-1.32,1.11) -- (-1.02,1.0);
\end{tikzpicture}

{\small الكرات الواحدية لمعايير $p$ في المستوي، متداخلةً حين
ينمو $p$ من $1$ (المعيّن) مرورًا بالقيمة $2$ (القرص) و $4$ (شبه
الإهليلج) وصولًا إلى $\infty$ (المربّع). وتحدّب كل كرة \emph{هو}
متراجحة منكوفسكي؛ والزوايا عند $p = 1$ و $p = \infty$ هي حيث
ينحلّ دفعةً واحدة كلٌّ من التحدّب التام ووحدانية أفضل التقريبات
وحالات التساوي في \cref{exo:b3:lp:12}.}
\end{omfigure}

\begin{proposition}[متسلسلة نويمان]\label{prop:b3:banach:neumann}
ليكن $E$ فضاء باناخ وليكن $T \in \mathcal L(E) = \mathcal L(E,E)$ مع $\vertiii T < 1$. عندئذٍ
يكون $I - T$ قابلًا للقلب في $\mathcal
L(E)$، مع $(I - T)^{-1} = \sum_{n\geq0}T^n$ (وهي متقاربة
\omterm{def:b3:banach:operator}{بمعيار المؤثر}). ومن ثَمّ فإن مجموعة المؤثرات القابلة للقلب
مفتوحة، والقلب متصل عليها.\index{متسلسلة نويمان}
\end{proposition}

\begin{proof}
تتقارب المتسلسلة تقاربًا مطلقًا ($\vertiii{T^n} \leq \vertiii
T^n$، وهي هندسية) في فضاء
باناخ $\mathcal L(E)$ (\cref{prop:b3:banach:llcomplete}؛ \cref{exo:b3:complete:1}(b)). وبالاختصار
التلسكوبي، $(I - T)\sum_{n \leq
N}T^n = I - T^{N+1} \to I$، وبالمثل من الجهة الأخرى. وأما الانفتاح: فإذا
كان $S$ قابلًا للقلب وكان $\vertiii{H} <
1/\vertiii{S^{-1}}$، فإن $S + H = S(I + S^{-1}H)$ مع $\vertiii{S^{-1}H} < 1$: أي إنه
قابل للقلب. وأما اتصال القلب: فيعطي تعبير المتسلسلة أن $\vertiii{(S+H)^{-1} - S^{-1}} =
O(\vertiii H)$
محليًا.
\end{proof}

\begin{example}[معادلة فولتيرا، محلولةً بنويمان]
\label{ex:b3:banach:volterraexample}
على $E = \mathcal C(\intcc01)$ لننظر في المعادلة التكاملية
\[
u(x) = 1 + \lambda\int_0^xu(t)\,\dd t,
\qquad\text{أي}\qquad u = \mathbf 1 + \lambda Tu,
\quad (Tu)(x) = \int_0^xu .
\]
وهنا $\vertiii T \leq 1$، ومنه فمن أجل $\abs\lambda < 1$ تنطبق \omterm{prop:b3:banach:neumann}{متسلسلة نويمان} مباشرةً:
$u = (I - \lambda T)^{-1}\mathbf 1 =
\sum_n\lambda^nT^n\mathbf 1$. وبالحساب، $T^n\mathbf 1 =
\frac{x^n}{n!}$، ومنه
\[
u(x) = \sum_{n\geq0}\frac{(\lambda x)^n}{n!} =
\eu^{\lambda x} ,
\]
كما يؤكّد الاشتقاق. بل أفضل: $\vertiii{T^n} \leq
\frac1{n!}$ (لأن النواة المتكررة تتقلص
عامليًا)، ومنه فإن $\sum\lambda^nT^n$ تتقارب من أجل \emph{كل} $\lambda$ ---
أي إن المؤثر $I - \lambda T$ قابل للقلب من أجل كل $\lambda
\in \C$، حتى
وإن كان $\vertiii{\lambda T} \geq 1$ في نهاية المطاف: فالمهم هو التناقص الطيفي للقوى،
لا المعيار الأول. ولمؤثرات فولتيرا هذا التناقص العاملي مضمَّنًا
فيها (وقد استثمر \cref{exo:b3:complete:4}(a) ذلك بالضبط)، ولهذا لا تعاني
مسائل القيم الابتدائية أبدًا من ظواهر الرنين التي تعانيها مسائل
القيم الحدّية (\cref{ch:b3:spectral}).
\end{example}

\section{هان--باناخ}

\begin{theorem}[هان--باناخ، الصيغة التحليلية]
\label{thm:b3:banach:hahnbanach}
ليكن $E$ فضاءً متجهيًا \emph{حقيقيًا}، وليكن $p \colon E \to \R$ \emph{تحت
خطي} (أي $p(x + y) \leq p(x) + p(y)$ و $p(tx) =
tp(x)$ من أجل $t \geq 0$)، وليكن
$F \subseteq E$ فضاءً جزئيًا وليكن $f
\colon F \to \R$ خطيًا مع $f \leq p$
على $F$. عندئذٍ يمتد $f$ إلى تطبيق خطي $\tilde f \colon E \to \R$ يحقق $\tilde f \leq p$ على
$E$.\index{مبرهنة هان--باناخ}
\end{theorem}

\begin{proof}
\emph{التمديد بخطوة واحدة.} ليكن $x_0 \notin F$؛ نمدّد $f$ إلى
$F \oplus \R x_0$ باختيار $\alpha = \tilde f(x_0)$ اختيارًا صحيحًا: إذ نحتاج، من أجل كل
$y \in F$ و $t > 0$، إلى
\[
f(y) + t\alpha \leq p(y + tx_0)
\quad\text{و}\quad
f(y) - t\alpha \leq p(y - tx_0),
\]
وهذا يختزل بعد القسمة على $t$ (بالتحت خطية) إلى
\[
\sup_{v \in F}\ \bigl[f(v) - p(v - x_0)\bigr]
\;\leq\; \alpha \;\leq\;
\inf_{u \in F}\ \bigl[p(u + x_0) - f(u)\bigr].
\]
ويوجد مثل هذا $\alpha$ إذا وفقط إذا كان كل طرف أيسر $\leq$ كل
طرف أيمن: وبالفعل $f(v) + f(u) = f(u + v) \leq p(u + v) \leq
p(u + x_0) + p(v - x_0)$، أي $f(v) - p(v - x_0) \leq p(u +
x_0) - f(u)$.

\emph{زورن.} نرتّب تمديدات $f$ المهيمَن عليها بالمقدار $p$ (وهي أزواج:
فضاء جزئي، ومؤثر دالي) بالتمديد؛ ولكل سلسلة الاتحادُ حدًّا
أعلى؛ ويجب أن يكون العنصر الأعظمي معرَّفًا على $E$ كلها، وإلا
ناقض التمديدُ بخطوة واحدة الأعظمية.
\end{proof}

\begin{corollary}\label{cor:b3:banach:hbnormed}
ليكن $E$ فضاءً معياريًا (حيث $K = \R$ أو $\C$).
\begin{enumerate}
  \item كل $f \in F'$ (حيث $F$ فضاء جزئي) يمتد إلى $\tilde f
        \in E'$ مع
        $\norm{\tilde f}_{E'} = \norm f_{F'}$.
  \item من أجل كل $x \neq 0$ يوجد $f \in E'$ يحقق $\norm f
        = 1$
        و $f(x) = \norm x$. وعلى وجه الخصوص، يفصل $E'$ نقطَ
        $E$، و $\norm x = \sup_{\norm f \leq
        1}\abs{f(x)}$.
  \item من أجل فضاء جزئي مغلق $F$ ومن أجل $x \notin F$، يوجد
        $f
        \in E'$ ينعدم على $F$ ويحقق $f(x) = d(x, F)$
        و $\norm f \leq 1$.
\end{enumerate}
\end{corollary}

\begin{proof}
(1) الحالة الحقيقية: نطبّق \cref{thm:b3:banach:hahnbanach} مع $p(x)
= \norm f_{F'}\,\norm x$ (وهو تحت خطي)؛
ويحقق التمديد $\pm\tilde f(x) = \tilde f(\pm x) \leq p(x)$، ومنه $\norm{\tilde f} \leq \norm f$، والمتراجحة $\geq$ هي القصر.
الحالة العقدية: ليكن $u = \operatorname{Re}f$، وهو مؤثر دالي حقيقي يحقق $\abs u \leq \norm f\norm\cdot$؛
ولاحظ $f(x) = u(x) -
\iu\,u(\iu x)$ (بالتحقق على الجزأين الحقيقي والتخيلي: $\operatorname{Im}f(x) = -\operatorname{Re}f(\iu x)$).
نمدّد $u$ تمديدًا خطيًا حقيقيًا بالحد نفسه، ونضع $\tilde f(x) =
\tilde u(x) - \iu\tilde u(\iu x)$: فهو خطي على
$\C$ (بتحقّق مباشر على الضرب في $\iu$)، ويمدّد $f$؛ وأما
المعيار: فمن أجل $x$ معطى نكتب $\tilde f(x) = r\eu^{\iu\theta}$، ثم $\abs{\tilde f(x)}
= \tilde f(\eu^{-\iu\theta}x) = \tilde u(\eu^{-\iu\theta}x)
\leq \norm f\,\norm x$.

(2) على $F = Kx$ نعرّف $f(tx) = t\norm x$: ومعياره $1$ على $F$؛ ثم نمدّده
حسب (1). وأما صيغة الثنوية: فالمتراجحة $\leq$ واضحة، والمتراجحة
$\geq$ بهذا $f$.

(3) على $F \oplus Kx$ نعرّف $f(y + tx) = t\,d(x, F)$؛ وعندئذٍ من أجل $t \neq 0$
لدينا $\norm{y + tx} = \abs t\,\norm{x + y/t} \geq \abs
t\,d(x, F) = \abs{f(y + tx)}$: أي $\norm f \leq 1$ على الفضاء الجزئي؛ ثم نمدّده
حسب (1).
\end{proof}

\begin{remark}\label{rem:b3:banach:bidual}
حسب (2)، يكون التطبيق القانوني $J \colon E \to E''$، $J(x)(f) =
f(x)$، تقييسًا
(\cref{exo:b3:banach:10}): أي إن كل فضاء معياري يقع داخل ثنويّه المضاعف.
وتُسمّى الفضاءات التي يكون فيها $J$ شاملًا
\emph{انعكاسية}\index{فضاء انعكاسي}؛ وتُظهر مسألة نهاية الأسبوع
أن $\ell^p$ (حيث $1 < p < \infty$) انعكاسي بينما $\ell^1$ ليس
كذلك.
\end{remark}

\section{ثلاثية بير}

\begin{theorem}[باناخ--شتاينهاوس، الحد المنتظم]
\label{thm:b3:banach:banachsteinhaus}
ليكن $E$ فضاء \emph{باناخ}، وليكن $F$ معياريًا، ولتكن $(T_i)_{i\in
I} \subseteq \mathcal L(E, F)$
عائلة تحقق $\sup_i
\norm{T_ix} < \infty$ من أجل كل $x \in E$. عندئذٍ
$\sup_i
\vertiii{T_i} < \infty$.\index{مبرهنة باناخ--شتاينهاوس}
\end{theorem}

\begin{proof}
المجموعات $F_n = \{x : \sup_i\norm{T_ix} \leq n\}$ مغلقة (بوصفها تقاطعات لسوابق كرات مغلقة)
وتغطّي $E$. وتعطي مبرهنة بير (\cref{thm:b3:complete:baire}) عددًا $n_0$ وكرة
$B(x_0, r) \subseteq F_{n_0}$. ومن أجل $\norm z < r$: $\norm{T_iz}
\leq \norm{T_i(x_0 + z)} + \norm{T_ix_0} \leq 2n_0$، ومنه $\vertiii{T_i} \leq 2n_0/r$ من أجل كل
$i$.
\end{proof}

\begin{corollary}\label{cor:b3:banach:bslimits}
إذا كان $E$ فضاء باناخ وتقاربت $T_n \in \mathcal L(E,F)$ \emph{نقطةً نقطة} (أي
$T_nx \to Tx$ من أجل كل $x$)، فإن $\sup_n
\vertiii{T_n} < \infty$ و $T \in \mathcal L(E, F)$ و $\vertiii T \leq \liminf \vertiii{T_n}$.
\end{corollary}

\begin{proof}
المتتاليات المتقاربة محدودة: أي الحد نقطةً نقطة؛ وتحدّ مبرهنة
باناخ--شتاينهاوس المعاييرِ بعدد $M$ ما؛ عندئذٍ $\norm{Tx}
= \lim\norm{T_nx} \leq M\norm x$ (لأن $T$
خطي بوصفه نهاية نقطةً نقطة)، والحد الأدقّ بالانتقال إلى
$\liminf$ في $\norm{T_nx} \leq \vertiii{T_n}\norm x$.
\end{proof}

\begin{theorem}[متسلسلات فورييه المتباعدة]
\label{thm:b3:banach:fourierdiverge}
توجد دوال متصلة دورية بالدور $2\pi$ تتباعد متسلسلة فورييه لها
عند $0$: أي $\sup_N\abs{S_N(f)(0)} =
\infty$. بل إن هذه الدوال $f$ تشكّل جزءًا كثيفًا من
$\mathcal
C(S^1)$.\index{متسلسلة فورييه، تباعدها}
\end{theorem}

\begin{proof}
نعمل في $E = (\mathcal C(S^1), \norm\cdot_\infty)$، وهو فضاء
باناخ، بالمؤثّرات
$\Lambda_N(f) = S_N(f)(0) =
\frac1{2\pi}\int_{-\pi}^{\pi}f(t)\,D_N(t)\,\dd t$
(نواة ديريكليه، السنة الجامعية 2). وكل $\Lambda_N$ متصل مع
\[
\norm{\Lambda_N} = \frac1{2\pi}\int_{-\pi}^\pi\abs{D_N(t)}\,\dd
t \;=\; L_N .
\]
(فالمتراجحة $\leq$ واضحة؛ وأما $\geq$: فنأخذ $f$ متصلة مع
$\norm f_\infty
\leq 1$ تقرّب $\operatorname{sign}D_N$ --- فللإشارة عدد منتهٍ من القفزات؛ وتنعيم
كل قفزة على فترة طولها $\varepsilon$ يغيّر التكامل بمقدار
$O(N\varepsilon)$.)
وتؤول \emph{ثوابت لوبيغ} $L_N$ إلى اللانهاية:
\[
L_N = \frac1{2\pi}\int_{-\pi}^{\pi}
\frac{\abs{\sin\bigl((N{+}\tfrac12)t\bigr)}}{\abs{\sin(t/2)}}
\,\dd t
\geq \frac{2}{\pi}\int_0^{\pi}
\frac{\abs{\sin\bigl((N{+}\tfrac12)t\bigr)}}{t}\,\dd t
= \frac{2}{\pi}\int_0^{(N+\frac12)\pi}\frac{\abs{\sin
u}}{u}\,\dd u
\]
(باستعمال $\abs{\sin(t/2)} \leq t/2$ على $[0, \pi]$، ثم بالتعويض $u = (N + \tfrac12)t$). وبالقطع
إلى أقواس:
\[
\int_{(k-1)\pi}^{k\pi}\frac{\abs{\sin u}}u\,\dd u
\geq \frac1{k\pi}\int_{(k-1)\pi}^{k\pi}\abs{\sin u}\,\dd u =
\frac{2}{k\pi},
\qquad\text{ومنه}\qquad
L_N \geq \frac{4}{\pi^2}\sum_{k=1}^N\frac1k
\xrightarrow[N\to\infty]{} \infty .
\]
ولو كان لكل دالة متصلة $f$ الشرطُ $\sup_N
\abs{\Lambda_N(f)} < \infty$، لفرضت مبرهنة
باناخ--شتاينهاوس أن $\sup_N\norm{\Lambda_N} < \infty$: وهو تناقض. إذن توجد دالة $f$ ---
بل مجموعة غير \omterm{rem:b3:complete:meagre}{هزيلة} وكثيفة من الدوال $f$ (وهي متمم $\bigcup_M\{f: \sup_N\abs{\Lambda_Nf}\leq M\}$،
وهو اتحاد قابل للعدّ لمجموعات مغلقة ليس لها داخل حسب ما سبق مطبَّقًا
في أي كرة، ومنه فهو هزيل) --- تحقق $\sup_N\abs{S_Nf(0)} =
\infty$.
\end{proof}

\begin{theorem}[التطبيق المفتوح]\label{thm:b3:banach:openmapping}
ليكن $E, F$ فضاءي باناخ وليكن $T \in \mathcal L(E, F)$ \emph{شاملًا}. عندئذٍ
يكون $T$ مفتوحًا: أي $T(B_E(0,1)) \supseteq
B_F(0, c)$ من أجل $c > 0$ ما. ومن ثَمّ فإن
لكل مؤثر محدود تقابلي بين فضاءي باناخ مقلوبًا محدودًا.
\index{مبرهنة التطبيق المفتوح}
\end{theorem}

\begin{proof}
نكتب $B = B_E(0,1)$. تعطي الشمولية $F = \bigcup_n
\overline{T(nB)} = \bigcup_n n\,\overline{T(B)}$؛ وتعطي مبرهنة بير
(\cref{thm:b3:complete:baire}) داخلًا للمجموعة $\overline{T(B)}$: أي $B_F(y_0, 4c) \subseteq
\overline{T(B)}$ ما. ونعيد
التمركز عند $0$: فمن أجل $\norm y < 4c$، يكون كلٌّ من $y_0 + y$
و $y_0$ نهايةً لصور $Tu_k$ و $Tv_k$ مع $u_k, v_k \in B$، ومنه
$y = \lim T(u_k - v_k)$ مع $u_k - v_k
\in 2B$: أي $B_F(0, 4c) \subseteq \overline{T(2B)}$، أي $B_F(0, 2c) \subseteq \overline{T(B)}$.

\emph{إزالة \omterm{def:b3:topology:interior}{الغلق}} (وهنا يدخل تمام $E$): ليكن $\norm y < c$.
نختار $x_1 \in \frac12 B$ يحقق $\norm{y -
Tx_1} < c/2$ ($B_F(0,2c) \subseteq \overline{T(B)}$ محجَّمة بالعامل
$\frac12$)؛ وبالتتابع $x_k \in 2^{-k}B$ مع $\norm{y -
T(x_1 + \dots + x_k)} < c\,2^{-k}$. وتتقارب المتسلسلة
$\sum x_k$ تقاربًا مطلقًا في فضاء باناخ $E$ إلى $x \in B$
(بمعيار $<
\sum 2^{-k} = 1$)، ويكون $Tx = y$ بالاتصال: أي $B_F(0, c)
\subseteq T(B)$. وينتج
انفتاح $T$ على المفتوحات الكيفية بالانسحاب والتحجيم؛ وأما
النتيجة الفرعية، فانفتاح $T$ يعني أن $T^{-1}$ متصل.
\end{proof}

\begin{corollary}[المعايير المتكافئة]\label{cor:b3:banach:equivnorms}
إذا كان فضاء متجهي تامًّا من أجل معيارين قابلين للمقارنة
($\norm\cdot_a \leq C\norm\cdot_b$)، فالمعياران متكافئان.
\end{corollary}

\begin{proof}
التطابق $(E, \norm\cdot_b) \to (E, \norm\cdot_a)$ محدود وتقابلي بين فضاءي باناخ: ومنه فمقلوبه
محدود.
\end{proof}

\begin{theorem}[المنحني البياني المغلق]\label{thm:b3:banach:closedgraph}
ليكن $E, F$ فضاءي باناخ وليكن $T \colon E \to F$ خطيًا. إذا كان المنحني
البياني $\Gamma = \{(x, Tx)\}$ مغلقًا في $E \times F$ (أي إن $x_n \to x$
و $Tx_n \to y$ يستلزمان $y = Tx$)، فإن $T$ محدود.\index{مبرهنة
المنحني البياني المغلق}
\end{theorem}

\begin{proof}
الفضاء $E \times F$ مزوَّدًا بالمعيار $\norm{(x,y)} = \norm x + \norm y$ فضاءُ باناخ؛ و $\Gamma$،
بوصفه فضاءً جزئيًا مغلقًا، فضاءُ باناخ. والإسقاط $\pi_E\colon \Gamma \to E$ محدود
وتقابلي، ومنه فمقلوبه $x \mapsto (x, Tx)$ محدود (\cref{thm:b3:banach:openmapping}): أي $\norm{Tx} \leq \norm{(x,
Tx)} \leq C\norm x$.
\end{proof}

\begin{method}\label{met:b3:banach:trilogy}
متى نلجأ إلى أي مبرهنة. \emph{هان--باناخ}: لإنتاج مؤثر دالي
بسلوك مفروض (تنظيم متجهة، أو انعدام على فضاء جزئي، أو تمديد من
فضاء جزئي) --- ولا حاجة إلى التمام. \emph{باناخ--شتاينهاوس}:
لتحويل معلومة نقطةً نقطة إلى حدود منتظمة --- وعادةً لإظهار أن
عملية نهاية متصلة، أو (بعكس النقيض) لإثبات \emph{التباعد} من
أجل عنصر ما، كما في متسلسلات فورييه. \emph{التطبيق المفتوح /
المنحني البياني المغلق}: للحصول على الاتصال مجانًا من التقابلية
الجبرية أو من خاصية انغلاق للمنحني البياني --- والاستعمال
النموذجي: مقارنة معيارين تامّين، أو إثبات الاتصال التلقائي.
وتتطلب مبرهنات بير الثلاث جميعًا تمام \emph{المنطلق}؛ وثمة
أمثلة مضادّة فيما عدا ذلك (\cref{exo:b3:banach:7}).
\end{method}

\section{الفضاءات الثنوية، بصورة ملموسة}

\begin{theorem}\label{thm:b3:banach:duals}
تقييسيًا: $(c_0)' \cong \ell^1$ و $(\ell^1)' \cong
\ell^\infty$، عبر المزاوجة
$\langle x, y\rangle = \sum_n
x_ny_n$.\index{ثنويات فضاءات المتتاليات}
\end{theorem}

\begin{proof}
نبرهن على $(c_0)' \cong \ell^1$؛ والتعيين الثاني هو \cref{exo:b3:banach:5}. نُلحق
بالعنصر $y \in \ell^1$ التطبيقَ $\Lambda_y(x)
= \sum x_ny_n$ (حيث $x \in c_0$): وهو
متقارب تقاربًا مطلقًا مع $\abs{\Lambda_y(x)} \leq \norm x_\infty\norm y_1$، ومنه $\norm{\Lambda_y} \leq \norm y_1$. وبالعكس، ليكن
$\Lambda \in
(c_0)'$؛ نضع $y_n = \Lambda(e_n)$ (حيث $e_n$ المتتاليات الواحدية). ومن أجل
أي $N$، نختبر على $x^{(N)} = \sum_{n \leq N}
\operatorname{sign}(\overline{y_n})\,e_n \in c_0$ (ومعياره $\leq
1$؛ وفي الحالة العقدية
نستعمل عوامل أحادية المعيار $\bar y_n/\abs
{y_n}$): فنجد $\Lambda(x^{(N)}) = \sum_{n\leq N}\abs{y_n} \leq
\norm\Lambda$. إذن
$y \in \ell^1$ مع $\norm y_1 \leq
\norm\Lambda$. وأخيرًا $\Lambda = \Lambda_y$: فكلاهما يتوافق على
العناصر $e_n$، ومنه على المتتاليات المنتهية، وهي كثيفة في $c_0$
(بالبتر: $\norm{x - \sum_{n \leq N}x_ne_n}_\infty = \sup_{n > N}\abs{x_n}
\to 0$ وذلك بالضبط لأن $x_n \to 0$)؛ والمؤثرات
الدالية المتصلة المتوافقة على \omterm{def:b3:topology:interior}{مجموعة كثيفة} متساوية. والتقابل
خطي وتقابلي وتقييسي ($\norm{\Lambda_y} = \norm y_1$ من المتراجحتين).
\end{proof}

\section{تمارين}

\begin{exercise}[$\star$]\label{exo:b3:banach:1}
احسب معايير المؤثرات:
(a) الإزاحتان $S(x_1, x_2, \dots) = (0, x_1, x_2, \dots)$ و $S^*(x_1, x_2, \dots) = (x_2, x_3, \dots)$ على $\ell^2$؛
(b) مؤثر الضرب $M_a x = (a_nx_n)$ على $\ell^2$، من أجل $a \in \ell^\infty$؛
(c) المؤثر الدالي $\Lambda(f) = \int_0^{1/2}f -
\int_{1/2}^1f$ على $\mathcal C(\intcc01)$ --- برهن على $\norm\Lambda = 1$ وعلى
أن المعيار \emph{غير مبلوغ}.
\end{exercise}

\begin{exercise}[$\star$]\label{exo:b3:banach:2}
ليكن $E$ فضاء باناخ، وليكن $T \in \mathcal L(E)$ قابلًا للقلب، وليكن $S$
يحقق $\vertiii{S - T} < 1/\vertiii{T^{-1}}$. برهن على أن $S$ قابل للقلب وقدّر $\vertiii{S^{-1} - T^{-1}}$. تطبيق:
إذا كانت جملة خطية $Tx = b$ \omterm{def:b3:groups:derived}{قابلة للحل} مع $T$ قابل للقلب، فإن
اضطرابًا صغيرًا بما يكفي للمؤثر $T$ يُبقيها \omterm{def:b3:groups:derived}{قابلة للحل} بطريقة
وحيدة، مع حدّ كمّي على تغيّر الحل.
\end{exercise}

\begin{exercise}[$\star\star$]\label{exo:b3:banach:3}
(a) برهن على أن $\ell^1$ و $\ell^\infty$ و $c_0$ فضاءات باناخ،
وعلى أن $c_0$ هو \omterm{def:b3:topology:interior}{الغلق} في $\ell^\infty$ لفضاء المتتاليات
المنتهية.
(b) برهن على $\ell^p \subseteq \ell^q$ مع $\norm\cdot_q \leq
\norm\cdot_p$ من أجل $1 \leq p \leq q \leq \infty$، وعلى أن
الاحتواء تام.
\end{exercise}

\begin{exercise}[$\star\star$]\label{exo:b3:banach:4}
ليكن $F \subseteq E$ فضاءً جزئيًا مغلقًا وليكن $x \notin F$.
باستعمال \cref{cor:b3:banach:hbnormed}، برهن على صيغة الثنوية
\[
d(x, F) = \max\bigl\{\abs{f(x)} : f \in E',\ \norm f \leq 1,\
f\restriction_F = 0\bigr\}
\]
(ولاحظ: أنها \emph{قيمة عظمى}). واستنتج أن $F = \bigcap\{\ker f : f
\in E',\ f\restriction_F = 0\}$: أي إن
الفضاءات الجزئية المغلقة هي بالضبط تقاطعات نوى المؤثرات
الدالية.
\end{exercise}

\begin{exercise}[$\star\star$]\label{exo:b3:banach:5}
برهن على $(\ell^1)' \cong \ell^\infty$ تقييسيًا، متّبعًا مخطط \cref{thm:b3:banach:duals} (فالمتتاليات
المنتهية كثيفة في $\ell^1$). وأين تنهار الحجّة من أجل
$(\ell^\infty)'$؟
\end{exercise}

\begin{exercise}[$\star\star$]\label{exo:b3:banach:6}
ليكن $E, F, G$ فضاءات معيارية مع $E$ فضاء باناخ، وليكن $B \colon E \times
F \to G$
ثنائي الخطية \omterm{def:b3:topology:continuity}{ومتصلًا} في كل متغيّر على حدة. برهن على أن $B$ متصل
(بالمعنى المشترك): أي $\norm{B(x,y)} \leq
C\norm x\norm y$. \emph{(طبّق باناخ--شتاينهاوس
على العائلة $(B(\cdot, y))_{\norm y \leq 1}$.)}
\end{exercise}

\begin{exercise}[$\star\star$]\label{exo:b3:banach:7}
(a) على $E = \mathcal C(\intcc01)$، قارن $\norm\cdot_\infty$ مع $\norm\cdot_1$: فالتطابق
$(E, \norm\cdot_\infty) \to
(E, \norm\cdot_1)$ محدود وتقابلي لكن مقلوبه غير محدود. وأي فرضية من فرضيات
\cref{cor:b3:banach:equivnorms} تخفق؟
(b) أظهر تطبيقًا خطيًا غير متصل من فضاء جزئي كثيف من $\ell^2$
إلى $K$ (مثلًا على المتتاليات المنتهية)، واشرح لماذا لا يناقض
هذا مبرهنة المنحني البياني المغلق.
\end{exercise}

\begin{exercise}[$\star\star\star$]\label{exo:b3:banach:8}
(هيلنغر--تُبليتز) ليكن $T \colon \ell^2 \to \ell^2$ خطيًا (ومعرَّفًا في كل مكان)
و\emph{متناظرًا}: أي $\langle Tx,
y\rangle = \langle x, Ty\rangle$ من أجل كل $x, y$، حيث $\langle
x, y \rangle = \sum x_n\bar y_n$.
برهن على أن $T$ محدود. \emph{(بالمنحني البياني المغلق: إذا كان
$x_k \to x$ و $Tx_k \to z$، فاختبر على $y$ كيفي.)} والعبرة: أن
المؤثرات المتناظرة غير المحدودة --- وهي هاميلتونيات ميكانيكا
الكم --- لا يمكن أبدًا أن تُعرَّف على الفضاء كله.
\end{exercise}

\begin{exercise}[$\star\star\star$]\label{exo:b3:banach:9}
(مبرهنة بويا في التربيع العددي) من أجل كل $n$، لتكن $\Lambda_n(f)
= \sum_{i=0}^{n} w_{i,n}f(x_{i,n})$
قاعدة تربيع عددي على $\mathcal C(\intcc01)$ (حيث $x_{i,n} \in \intcc01$، $w_{i,n} \in
\R$). برهن على
أن $\Lambda_n(f) \to \int_0^1f$ من أجل \emph{كل} دالة متصلة $f$ إذا وفقط إذا: (أ)
كان $\Lambda_n(P) \to \int_0^1P$ من أجل كل كثير حدود $P$، و (ب) كان $\sup_n\sum_i\abs{w_{i,n}} <
\infty$.
\emph{(احسب $\norm{\Lambda_n}$؛ واستعمل باناخ--شتاينهاوس وفايرشتراس.)}
وتحقق من أن القواعد ذات الأوزان \emph{الموجبة} والمضبوطة على
الثوابت تحقق (ب) تلقائيًا.
\end{exercise}

\begin{exercise}[$\star\star$]\label{exo:b3:banach:10}
برهن على أن $J \colon E \to E''$ المعطى بالعلاقة $J(x)(f) = f(x)$ تقييسٌ خطي
(استعمل \cref{cor:b3:banach:hbnormed}(2))، وعلى أنه شامل حين $\dim E < \infty$. وبرهن أيضًا
على أنه إذا كان $E'$ قابلًا للفصل فكذلك $E$. \emph{(اختر $x_n$
تنظّم تقريبًا متتاليةً كثيفة من $E'$ وبرهن على أن الفضاء
المولَّد المغلق لها هو $E$، بواسطة \cref{cor:b3:banach:hbnormed}(3).)}
\end{exercise}

\begin{exercise}[$\star\star$]\label{exo:b3:banach:11}
(فضاءات القسمة) ليكن $E$ فضاء باناخ وليكن $F \subseteq
E$ فضاءً جزئيًا
\emph{مغلقًا}. نعرّف على $E/F$
\[
\norm{\bar x} \;=\; d(x, F) = \inf_{y\in F}\norm{x - y} .
\]
(a) برهن على أن هذا \emph{معيار} معرَّف تعريفًا سليمًا على $E/F$
(وأين يدخل انغلاق $F$؟)، وعلى أن الإسقاط $\pi
\colon E \to E/F$ يحقق $\vertiii\pi \leq 1$
ويرسل الكرة الواحدية المفتوحة على الكرة الواحدية المفتوحة.
(b) برهن على أن $E/F$ تام. \emph{(استعمل محك المتسلسلات في
\cref{exo:b3:complete:1}(b): إذا أُعطيت أصناف $\bar x_k$ تحقق $\sum\norm{\bar x_k} < \infty$، فارفع
كلًّا منها إلى $x_k \in E$ يحقق $\norm{x_k} \leq \norm{\bar x_k} + 2^{-k}$ واجمع في $E$.)}
(c) احسب: من أجل $E = c$ (المتتاليات المتقاربة) ومن أجل $F =
c_0$،
برهن على $c/c_0 \cong K$ تقييسيًا بواسطة $\bar x \mapsto
\lim_nx_n$.
\end{exercise}

\begin{exercise}[$\star\star$]\label{exo:b3:banach:12}
(الإسقاطات المحدودة والفضاءات الجزئية المتمَّمة) ليكن $E$ فضاء
باناخ وليكن $P \colon E \to E$ خطيًا مع $P^2 = P$ (أي إسقاط \omterm{def:b3:galois:algebraic}{جبري})،
و $V = \operatorname{im}P$، و $W =
\ker P$.
(a) لنفترض أن $P$ محدود. برهن على أن $V$ و $W$ مغلقان وأن
$E = V \oplus W$ بالتفكيك $x = Px + (x -
Px)$.
(b) وبالعكس، لنفترض $E = V \oplus W$ مع كون $V,
W$ مغلقين
\emph{كليهما}، وليكن $P$ الإسقاطَ على $V$ على امتداد $W$. برهن
على أن $P$ محدود. \emph{(بالمنحني البياني المغلق: إذا كان
$x_n \to x$ و $Px_n \to z$، فإن $z \in V$ و $x_n - Px_n \to x - z \in
W$، وتعيّن
وحدانية التفكيك $z =
Px$.)}
(c) استنتج \emph{التكافؤ}: أن فضاءً جزئيًا $V$ يقبل إسقاطًا
محدودًا إذا وفقط إذا كان مغلقًا وله متمّم \omterm{def:b3:galois:algebraic}{جبري} مغلق --- ولاحظ
(دون برهان) أن ثمة فضاءات جزئية مغلقة بلا هذه الخاصية (والمثال
الكلاسيكي هو $c_0$ داخل $\ell^\infty$): فإن فضاءات هيلبرت،
حيث يفي $V^\perp$ بالغرض دائمًا (\cref{ch:b3:hilbert})، هي الاستثناء لا
القاعدة.
\end{exercise}

\section{مسألة: ثنوية الفضاءات \texorpdfstring{$\ell^p$}%
{lp}}

\begin{problem}[{مسألة نهاية الأسبوع --- $(\ell^p)' = \ell^q$، والانعكاسية،
وغرابة $\ell^\infty$}]
\label{pb:b3:banach:1}
لنثبّت $1 < p < \infty$ وليكن $q$ الأُسَّ المرافق، أي $\frac1p + \frac1q = 1$.
والمزاوجة في كل ما يلي هي $\langle x,
y\rangle = \sum_n x_ny_n$.

\medskip
\noindent\textbf{الجزء الأول --- هولدر ومنكوفسكي من أجل
المتتاليات.}
\begin{enumerate}
  \item (متراجحة يونغ) من أجل $a, b \geq 0$، برهن على $ab \leq
        \frac{a^p}p + \frac{b^q}q$،
        باستعمال تقعّر $\log$ أو بدراسة $t \mapsto \frac{t^p}p + \frac1q - t$.
  \item (هولدر) استنتج: $\abs{\langle x, y\rangle} \leq
        \norm x_p\norm y_q$ من أجل $x \in \ell^p$
        و $y \in \ell^q$؛ وعيّن حالة التساوي.
  \item (منكوفسكي) استنتج متراجحة المثلث من أجل $\norm\cdot_p$.
        \emph{(اكتب $\abs{x_n + y_n}^p \leq
        \abs{x_n}\,\abs{x_n{+}y_n}^{p-1} +
        \abs{y_n}\,\abs{x_n{+}y_n}^{p-1}$ وطبّق هولدر على كل حد.)}
  \item برهن على أن $\ell^p$ تام وعلى أن المتتاليات المنتهية
        كثيفة فيه.
\end{enumerate}

\medskip
\noindent\textbf{الجزء الثاني --- الثنوية $(\ell^p)' =
\ell^q$.}
\begin{enumerate}[resume]
  \item من أجل $y \in \ell^q$، برهن على أن $\Lambda_y(x) = \langle
        x, y\rangle$ يعرّف $\Lambda_y \in (\ell^p)'$
        مع $\norm{\Lambda_y} \leq \norm y_q$، وبالاختبار على
        $x_n = \abs{y_n}^{q-1}\operatorname{sign}(y_n)$
        (مبتورًا ومنظَّمًا على النحو الملائم)، على أن $\norm{\Lambda_y} = \norm y_q$.
  \item وبالعكس، إذا أُعطي $\Lambda \in (\ell^p)'$، نضع $y_n =
        \Lambda(e_n)$؛ برهن على
        $y \in \ell^q$ مع $\norm y_q \leq
        \norm\Lambda$ \emph{(بالاختبار على المبتورات
        كما في السؤال 5 وبجعل طول البتر ينمو)}، واخلص إلى
        $\Lambda = \Lambda_y$: أي إن التطبيق $y \mapsto \Lambda_y$ تماثل تقييسي $\ell^q \to (\ell^p)'$.
  \item استنتج أن $\ell^p$ \emph{انعكاسي} من أجل $1 < p <
        \infty$:
        فبتركيب الثنويتين، يأتي كل عنصر من $(\ell^p)''$ من
        $\ell^p$؛ وتحقق بعناية من أن المركّب هو $J$ القانوني.
\end{enumerate}

\medskip
\noindent\textbf{الجزء الثالث --- $\ell^1$ و $\ell^\infty$
حيوانان مختلفان.}
\begin{enumerate}[resume]
  \item برهن على أن $\ell^p$ (حيث $1 \leq p < \infty$) و $c_0$ قابلان
        للفصل، لكن $\ell^\infty$ ليس كذلك. \emph{(فمتتاليات
        الدلالة لأجزاء $\N$، وعددها غير قابل للعدّ، تبعد كلٌّ منها
        عن الأخرى مسافة $1$.)}
  \item استنتج من \cref{exo:b3:banach:10} أن $(\ell^1)'
        \cong \ell^\infty$ لكن $(\ell^\infty)' \not\cong
        \ell^1$: أي إن
        $\ell^1$ \emph{ليس} انعكاسيًا. \emph{(إذ لو كان
        $(\ell^\infty)'$ هو $\ell^1$، لكان قابلًا للفصل، فيفرض
        قابلية $\ell^\infty$ للفصل.)}
  \item (نهاية باناخ، صراحةً) على $\ell^\infty_\R$، لتكن
        $p(x) = \limsup_n \frac{x_1 + \dots + x_n}{n}$. برهن على أن $p$ تحت خطي، وعلى أنه على الفضاء
        الجزئي $c$ للمتتاليات المتقاربة يحقق $\mathrm{LIM}(x) = \lim x$ الشرطَ
        $\mathrm{LIM} \leq p$. ومدّده بهان--باناخ إلى $\mathrm{LIM} \colon \ell^\infty_\R \to \R$ وبرهن على أن
        $\mathrm{LIM}$ موجب ($x \geq 0 \Rightarrow
        \mathrm{LIM}(x) \geq 0$)، وصامد بالإزاحة ($\mathrm{LIM}(x_2, x_3, \dots) = \mathrm{LIM}(x)$)،
        ويمدّد النهاية، ويحقق $\liminf x \leq
        \mathrm{LIM}(x)\leq \limsup x$.
  \item برهن على أن مثل هذا $\mathrm{LIM}$، منظورًا إليه في
        $(\ell^\infty)'$، \emph{ليس} من الشكل $\Lambda_y$ من
        أجل أي $y \in \ell^1$؛ واخلص من جديد إلى $(\ell^\infty)'
        \neq \ell^1$.
        \emph{(قيّمه على المتتاليات الواحدية $e_n$، ثم على
        المتتالية الثابتة $1$.)}
  \item قيّم $\mathrm{LIM}$ على $(0,1,0,1,\dots)$، وبرهن على أنه لا يمكن
        وجود تمديد \emph{ضربي} صامد بالإزاحة للنهاية \emph{(انظر
        في $x =
        (0,1,0,1,\dots)$ وفي $x\cdot Sx$ حيث $S$ الإزاحة)}.
\end{enumerate}

\medskip
\noindent\textbf{الجزء الرابع --- خاتمة: لماذا تهمّ
\omterm{rem:b3:banach:bidual}{الانعكاسية}.}
\begin{enumerate}[resume]
  \item باستعمال \cref{cor:b3:banach:bslimits} والسؤال 6، برهن على أن لكل متتالية
        محدودة من $\ell^p$ (حيث $1 < p <
        \infty$) متتاليةً جزئية
        $(x^{(k)})$ تتقارب \emph{ضعيفًا}: أي إن $\Lambda(x^{(k)})$ تتقارب
        من أجل كل $\Lambda \in (\ell^p)'$. \emph{(بالاستخراج القطري على
        الإحداثيات القابلة للعدّ؛ وعيّن النهاية الضعيفة في $\ell^p$
        باستعمال الحد المنتظم للمعايير وهولدر.)} وبرهن بمثال
        (المتتالية $e_n$ في $\ell^1$، مقابل عناصر مختارة بعناية
        من $\ell^\infty$) على أن ذلك يخفق في $\ell^1$: فالتراص
        الضعيف امتياز للفضاءات \omterm{rem:b3:banach:bidual}{الانعكاسية}.
\end{enumerate}

\medskip
\noindent\textbf{الجزء الخامس --- \omterm{def:b3:topology:topology}{الطوبولوجيا} الضعيفة في
العمل، ومفاجأة شور.} نكتب $x^{(k)} \rightharpoonup x$ في فضاء معياري $E$
(\emph{التقارب الضعيف}) حين يكون $\Lambda(x^{(k)}) \to \Lambda(x)$ من أجل كل
$\Lambda \in E'$.
\begin{enumerate}[resume]
  \item أكمل الإحصاء: برهن على $(c_0)' \cong \ell^1$ تقييسيًا، بمخطط
        السؤالين 5 و 6 (وما الذي يحل محل متتاليات الاختبار؟).
        وكوِّن السلسلة $c_0 \to \ell^1 \to \ell^\infty \to \dots$ للثنويات المتتالية وحدّد أين
        تخفق \omterm{rem:b3:banach:bidual}{الانعكاسية}.
  \item برهن على أن كل متتالية متقاربة ضعيفًا في فضاء باناخ
        محدودة: انظر إلى العناصر $x^{(k)}$ عبر الغمر القانوني
        $J$ بوصفها مؤثرات دالية على $E'$ وطبّق
        باناخ--شتاينهاوس (\cref{thm:b3:banach:banachsteinhaus}) --- وعلى أي فضاء باناخ،
        ولماذا يكون التمام متاحًا هناك؟
  \item برهن على أنه في $\ell^p$ مع $1 < p < \infty$: يكون
        $x^{(k)}
        \rightharpoonup x$ إذا وفقط إذا كان $\sup_k\norm{x^{(k)}}_p <
        \infty$ و $x^{(k)}_n \to x_n$ من أجل كل
        إحداثي $n$ \emph{(ويستعمل أحد الاتجاهين
        باناخ--شتاينهاوس عبر الغمر القانوني؛ ومن أجل الآخر،
        قرّب $y
        \in \ell^q$ بمتتاليات منتهية)}. واستنتج $e_k
        \rightharpoonup 0$ في
        $\ell^2$ بينما $\norm{e_k}_2 =
        1$: أي إن النهايات الضعيفة يمكن أن
        تفقد كتلة.
  \item برهن على أن المعيار نصف متصل من الأسفل ضعيفًا: أي إن
        $x^{(k)} \rightharpoonup x$ يستلزم $\norm x \leq
        \liminf_k\,\norm{x^{(k)}}$ \emph{(اختر مؤثرًا داليًا منظّمًا
        للعنصر $x$، \cref{cor:b3:banach:hbnormed})}.
  \item (رادون--ريس في $\ell^2$) برهن على أنه في $\ell^2$، يستلزم
        التقارب الضعيف مع تقارب المعايير التقاربَ بالمعيار
        \emph{(انشر $\norm{x^{(k)} -
        x}_2^2$)}. وأعطِ مثالًا مضادًّا للعبارة
        نفسها دون فرضية المعيار.
  \item (شور، الخطوة 1) ليكن $x^{(k)} \rightharpoonup 0$ في $\ell^1$ ولنفترض جدلًا
        أن $\norm{x^{(k)}}_1 \geq \delta > 0$ على متتالية جزئية. برهن أولًا على $x^{(k)}_n \to 0$ من
        أجل كل $n$ (وأي المؤثرات الدالية؟)، ثم أنشئ بالتتابع
        أدلّة $k_1 < k_2 < \cdots$ وأعدادًا صحيحة $0 = N_0 < N_1 < N_2 < \cdots$ بحيث تتركّز كتلة
        $x^{(k_j)}$ على الكتلة $B_j =
        \intoc{N_{j-1}}{N_j}$:
        \[
        \sum_{n \in B_j}\bigl|x^{(k_j)}_n\bigr| \geq
        \norm{x^{(k_j)}}_1 - \frac\delta{10} .
        \]
  \item (شور، الخطوة 2) نعرّف $y \in \ell^\infty$ بالعلاقة $y_n =
        \operatorname{sign}\bigl(x^{(k_j)}_n\bigr)$ من أجل
        $n \in
        B_j$. برهن على $\langle x^{(k_j)}, y\rangle \geq
        \norm{x^{(k_j)}}_1 - \frac{2\delta}{10} \geq
        \frac{8\delta}{10}$ واستخرج تناقضًا مع $x^{(k)} \rightharpoonup 0$.
        واخلص إلى \emph{مبرهنة شور}: أي إن المتتاليات المتقاربة
        ضعيفًا في $\ell^1$ تتقارب بالمعيار.
  \item استنتج أنه ليست للمتتالية $(e_k)$ متتالية جزئية متقاربة ضعيفًا
        في $\ell^1$ (لأن نهايتها المرشّحة الوحيدة هي $0$
        إحداثيًا --- ثم نطبّق شور)، فتستعيد إخفاق التراص الضعيف
        في السؤال 13؛ وحلّ المفارقة الظاهرية: ففي $\ell^1$
        يتطابق التقارب الضعيف والتقارب بالمعيار
        \emph{للمتتاليات}، ومع ذلك تختلف \emph{الطوبولوجيتان}
        الضعيفة والمعيارية وتبقى المجموعات المحدودة غير متراصة
        ضعيفًا بالمتتاليات --- ولا تناقض في ذلك، بل مجرّد إخفاق
        \omterm{rem:b3:banach:bidual}{الانعكاسية}.
  \item (جدول تركيبي) من أجل $E \in \{c_0,\ \ell^1,\ \ell^p\
        (1{<}p{<}\infty),\ \ell^\infty\}$، كوِّن جدولًا يضم: \omterm{def:b3:banach:operator}{الفضاء
        الثنوي}؛ والقابلية للفصل؛ \omterm{rem:b3:banach:bidual}{والانعكاسية}؛ وهل تقبل
        المتتاليات المحدودة متتاليات جزئية متقاربة ضعيفًا؛
        وخاصية مميّزة واحدة لكل فضاء، مبرَّرة بسطر واحد من هذه
        المسألة.
\end{enumerate}

\medskip
\noindent\textbf{الجزء السادس --- تكملات: أقرب النقط، والتقارب
بالمتوسطات، وقيمة نهاية باناخ.}
\begin{enumerate}[resume]
  \item (أقرب النقط: عائد \omterm{rem:b3:banach:bidual}{الانعكاسية}) ليكن $F$ فضاءً جزئيًا
        مغلقًا من $\ell^p$ (حيث $1 < p < \infty$) وليكن
        $x \in \ell^p$. برهن على أن $d = \operatorname{dist}(x,
        F)$ \emph{مبلوغة}:
        استخرج من متتالية مصغّرة متتاليةً جزئية متقاربة ضعيفًا
        (السؤال 13)، وأبقِ النهاية الضعيفة داخل $F$ ببناء مؤثر
        دالي، بواسطة هان--باناخ، ينعدم على $F$ ولا ينعدم عند
        نقطة خارجه، ثم اخلص بالسؤال 17. ثم برهن على أن هذا
        الامتياز ليس شاملًا: ففي $c_0$، من أجل $\Lambda(x) = \sum_n2^{-n}x_n$، برهن
        على أن $\norm\Lambda = 1$ غير مبلوغة على الكرة الواحدية، وأثبت
        صيغة المسافة $\operatorname{dist}(x, \ker\Lambda) =
        \abs{\Lambda(x)}$، واستنتج أنه ليس لأي $x \notin
        \ker\Lambda$ أقربُ
        نقطة في الفضاء الفائق المغلق $\ker\Lambda$.
  \item (باناخ--ساكس في $\ell^2$) ليكن $x^{(k)}
        \rightharpoonup 0$ في $\ell^2_{\R}$
        مع $\norm{x^{(k)}}_2 \leq C$. أنشئ متتالية جزئية $(y_j)$ تحقق $\abs{\langle y_i, y_j\rangle} \leq
        \frac1j$ من
        أجل كل $i < j$، واستنتج
        \[
        \Bigl\lVert\frac{y_1 + \dots + y_m}m\Bigr\rVert_2^2
        \leq \frac{C^2 + 2}m \longrightarrow 0 :
        \]
        أي إن متوسطات تشيزارو تتقارب \emph{بالمعيار} بعد
        الاستخراج. وتحقق على $(e_k)$، التي معيار متوسطاتها
        $\frac1{\sqrt m}$: فالتقارب الضعيف، عديم الجدوى للمتتالية نفسها
        (السؤال 16)، يصير تقاربًا بالمعيار للمتوسطات.
  \item (قيمة نهاية باناخ) لتكن $L$ أي نهاية باناخ (السؤال 10)
        ولتكن $A_mx = \frac1m(x + Sx + \dots
        + S^{m-1}x)$. برهن على $L(A_mx) = L(x)$ وعلى $\liminf A_mx
        \leq L(x) \leq \limsup A_mx$؛
        واستنتج أن \emph{جميع} نهايات باناخ تتوافق على
        المتتاليات الدورية، بالقيمة المساوية للمتوسط على دور ---
        أي $\frac13$ على $(1, 0, 0, 1,
        0, 0, \dots)$، وهو متسق مع القيمة $\frac12$
        في السؤال 12. ثم برهن على أن التوافق يخفق في العموم:
        فمن أجل المتتالية الكتلية $x$ المساوية $1$ على $\intoc{3^{j-1}}{3^j}$
        من أجل $j$ زوجي و $0$ فيما عدا ذلك، برهن على أن متوسطات
        تشيزارو تتذبذب بين $\leq
        \frac13$ و $\geq \frac23$، وابنِ
        نهايتَي باناخ $L_\pm$ تحققان $L_-(x) \leq \frac13 < \frac23
        \leq L_+(x)$ \emph{(بالتمديد من
        $c \oplus \R x$ بالقيمتين القصويين المقبولتين $\pm$:
        وتحقق من أن $\Lambda(y + tx) = \lim y + t\,p(x)$ مهيمَن عليه بالمقدار التحت خطي $p$
        في السؤال 10)}.
\end{enumerate}
\end{problem}
